\chapter{ماذا تحسب عصبونات \nt{SERO}}
\label{sec:serocomputer}

\section{أول قناة بثّ}

حتى الآن تخاطبت كل عصبونات الكتاب عبر ناقل العناوين: كان \nt{GLUT} و\nt{GABA} يرسلان النبضات إلى متلقّين محددين، وعرفنا ماذا تحسب شبكة كهذه وبأي لغة تتكلم. ننتقل الآن إلى الناقل الثاني من القسم~\ref{sec:buses}، وهو ناقل البثّ. أولى قنواته \nt{SERO}. وكما من قبل، لا نسمح لأنفسنا إلا بما يوجد في الكائن الحي.

في هذا الفصل تظهر ثلاث أدوات ستستعملها بعد ذلك كل قنوات البثّ: عصبون يعمل من تلقاء نفسه ما دام يصله دفق خلفي ثابت، إلى أن يوقفه شيء؛ وسلك هو في الحقيقة حجم من النسيج؛ وتركيز بطيء بدل النبضات المنفردة. وستُبنى \nt{DOPA} و\nt{NORA} و\nt{ACET} في الفصول~\babelsublr{\ref{sec:dofa}--\ref{sec:acet}} من القطع نفسها.

من وجهة نظر المهندس، يختلف عصبون \nt{SERO} في أربعة أمور مهمة.

\emph{الأول: المتلقي هو الذي يختار الإشارة.} للسيروتونين مستقبلات بالإشارتين، والقاعدة~\eqref{eq:dalesign} لا تنطبق هنا: لا يحدد الإشارةَ المرسِلُ، بل مستقبِلُ المتلقي (القضية~\ref{prop:receptor})~\cite{BarnesSharp1999}.

\emph{الثاني: عصبون \nt{SERO} يعمل من تلقاء نفسه إذا توفّر دفق خلفي ثابت.} في اليقظة يطلق النبضات بانتظام، بضع مرات في الثانية، ولا يحتاج إلى دفعة لكل نبضة: إنه ناظمة إيقاع. لكنه يحتاج إلى دفق خلفي دائم. في شريحة الدماغ، حيث لا يوجد هذا الدفق، تصمت معظم هذه الخلايا؛ ويعود الإيقاع المنتظم إذا أُعطي النورأدرينالين عبر مستقبلات $\alpha_1$، ويخمده حجبُها~\cite{VandermaelenAghajanian1983}. في الكائن الحي يأتي هذا الدفق الخلفي من \nt{NORA}. وبلغة الدارات هذا مذبذب يحتاج إلى تغذية كهربائية: ما دامت موجودة يعمل وحده، إلى أن يوقفه التثبيط.

\emph{الثالث: إنه بطيء.} تكاد كل مستقبلات السيروتونين تكون أيضية التوجه: لا تفتح قناة بنفسها، بل تطلق داخل الخلية سلسلة من التفاعلات. والاستجابة بطيئة: يرتفع التيار المثبِّط عبر المستقبل الرئيسي ثم يهبط في غضون ثانية تقريبًا~\cite{CourtneyFord2016}، أي أطول بعشرات المرات من استجابة مشبك \nt{GLUT} السريع.

\emph{الرابع: يحرّر الناقل لا في سلك بل في حجم.} في المناطق التي تصل إليها مخارج عصبونات \nt{SERO} كثيرًا ما يخرج السيروتونين لا إلى شق ضيق، بل إلى الحيز بين الخلايا، وينتشر حوله~\cite{Bunin1998}. يحسّ المتلقي بالتركيز، ولا يستطيع أن يعرف من أي مرسِل جاءت الجزيئة.

مع هذه الأزمنة لم تعد النبضات المنفردة مهمة، لذلك سنصف الشبكة بالترددات $r_j$ وبالتركيزات. يزداد تركيز السيروتونين $c$ في الحجم الذي تحرّر فيه العصبونات $j$ الناقل بفعل التحرير، وينقص لأن الناقل يُضَخّ عائدًا إلى الخلايا:
\begin{empheq}[box=\keybox]{equation}
\tau_c\,\frac{\dd c}{\dd t} \;=\; -c \;+\; \kappa\sum_j r_j ,
\qquad
r_i \;=\; f_i\bigl(b_i + s_i\,g\,c\bigr), \quad s_i=\pm1 .
\label{eq:seroneuron}
\end{empheq}
هذه طريقة أخرى لقراءة سيل النبضات تكمّل القضية~\ref{prop:reader}: الحجم لا يرى النبضات المنفردة ولا ترتيبها، بل التردد فقط، متوسَّطًا على الزمن $\tau_c$. المعادلة الأولى هي دارة \foreignlanguage{english}{RC} نفسها التي رأيناها في الغشاء: $\tau_c$ عمر الناقل في الحجم، ولم يُقَس مباشرة، فنفترضه من رتبة الثانية؛ و$\kappa$ مقدار الناقل الذي تعطيه نبضة واحدة. والمعادلة الثانية تصف المتلقي: $b_i$ دفقه الخاص، وهو موجب عند ناظمة الإيقاع (ويدخل فيه الدخل الخلفي الثابت من \nt{NORA})؛ و$g$ حساسية المستقبل؛ والإشارة $s_i$ يختارها المتلقي؛ و$f_i$ خاصية نقل غير متناقصة.

\section{بوابة \foreignlanguage{english}{NOT} من عصبون واحد}

لنأخذ ناظمة إيقاع ذات مستقبل مثبِّط، $s=-1$. ما دام لا سيروتونين حولها فهي تعمل على دفقها الخاص $b$. وعندما يظهر السيروتونين ينقص الدفق بمقدار $gc$، وإذا كان $gc>b$ يصمت العصبون. هذه بوابة \foreignlanguage{english}{NOT}: يوجد خرج حين لا يوجد دخل. وقد كلّفتنا عصبونًا واحدًا (الشكل~\ref{fig:serogates}). والقضية~\ref{prop:monotone} لا تنطبق هنا: فقد اشترطت أوزانًا موجبة.

إذا أثّر في ناظمة الإيقاع نفسها سيروتونين من حجمين مختلفين، فإنها تصمت متى وُجد السيروتونين في أحدهما على الأقل. هذه بوابة \foreignlanguage{english}{NOR}، ومن بوابات \foreignlanguage{english}{NOR} يمكن بناء أي دالة منطقية. ولا حاجة هنا إلى الترميز ثنائي الأسلاك الذي تستعمله شبكة \nt{GLUT}: غياب الإشارة تحسبه عصبونات تعمل إلى أن يوقفها شيء.

\begin{figure}[t]
\centering
\begin{tikzpicture}[>=Stealth,
  nrn/.style={circle,draw,very thick,fill=blue!6,minimum size=0.95cm,inner sep=0pt},
  pace/.style={nrn,double,double distance=1.2pt},
  inh/.style={-{Bar[width=7pt]},thick,red!70!black}]
\node[pace] (N1) at (0,0) {};
\draw[inh] (-1.6,0) node[left,black] {$c$} -- (N1);
\draw[->,very thick,red!70!black] (N1) -- (1.3,0) node[right,black] {$\bar c$};
\node[below=0.55cm] at (0,0) {\small \foreignlanguage{english}{NOT}};
\node[pace] (N2) at (5,0) {};
\draw[inh] (3.4,0.45) node[left,black] {$c_1$} -- (N2);
\draw[inh] (3.4,-0.45) node[left,black] {$c_2$} -- (N2);
\draw[->,very thick,red!70!black] (N2) -- (6.3,0) node[right,black] {$\overline{c_1\vee c_2}$};
\node[below=0.55cm] at (5,0) {\small \foreignlanguage{english}{NOR}};
\end{tikzpicture}
\caption{المنطق بعصبونات \nt{SERO}. الدائرة المزدوجة ناظمة إيقاع: مع دفق خلفي ثابت تعمل وحدها إلى أن تُثبَّط. الخط المنتهي بشرطة مستقبل مثبِّط، $s=-1$؛ و$c$ و$c_1$ و$c_2$ تركيزات السيروتونين في الأحجام المحيطة بالمتلقي. على اليسار \foreignlanguage{english}{NOT}، وعلى اليمين \foreignlanguage{english}{NOR}.}
\label{fig:serogates}
\end{figure}

\begin{keyblock}
\begin{proposition}[اكتمال منطق \nt{SERO}]
\label{prop:serocomplete}
إذا كانت في الشبكة نواظم إيقاع ذات مستقبلات مثبِّطة، ولم تختلط التحريرات في الأحجام المختلفة، استطاعت الشبكة~\eqref{eq:seroneuron} أن تحسب أي دالة منطقية، ويُنقل كل بت بسلك واحد.
\end{proposition}
\end{keyblock}

من الناحية المنطقية شبكة \nt{SERO} أقوى من \nt{GLUT}، لكن الشرط ”لا تختلط التحريرات في الأحجام المختلفة“ لا يتحقق في الدماغ (القسم~\ref{sec:volume}).

\section{التغذية الراجعة السالبة}

لعصبونات السيروتونين مستقبلات مثبِّطة عليها هي نفسها~\cite{Sprouse1987}. فالسيروتونين الذي حرّرته يعود إليها ويكبحها. هذه دارة مألوفة للمهندس: مضخّم بتغذية راجعة سالبة (الشكل~\ref{fig:serofeedback}).

ما المقيس هنا وما النموذج؟ في نواة الرفاء نفسها، حيث تقع عصبونات \nt{SERO}، يثبّط السيروتونين الجيران لا عبر حجم مشترك، بل من نقطة إلى نقطة عبر المشابك: يزيله بروتين النقل بسرعة ولا يدعه يتراكم خارج الشق. والتحرير المنفرد لا يعطي إلا توقفًا قصيرًا في التفريغ، أما التردد المتوسط فلا يتغير~\cite{CourtneyFord2016}. لذلك فالحلقة أدناه، المغلقة عبر حجم مشترك، نموذج: نحسب ما كانت ستفعله تغذية راجعة كهذه لو عملت على مستوى الترددات المتوسطة.

\begin{figure}[t]
\centering
\begin{tikzpicture}[>=Stealth,
  nrn/.style={circle,draw,very thick,fill=blue!6,minimum size=0.95cm,inner sep=0pt},
  pace/.style={nrn,double,double distance=1.2pt},
  blk/.style={draw,very thick,rounded corners=2pt,fill=orange!10,minimum height=0.9cm,minimum width=2.2cm,align=center,font=\small},
  inh/.style={-{Bar[width=7pt]},thick,red!70!black}]
\node[pace] (N) at (0,0) {};
\node[blk] (V) at (3.6,0) {الحجم\\$\tau_c$};
\draw[->,very thick] (N) -- node[above] {\small $r$} (V);
\draw[->,very thick] (V) -- node[above] {\small $c$} (6.0,0) node[right] {\small إلى كل المتلقين};
\draw[inh] (V.south) -- ++(0,-0.9) -| node[pos=0.25,below] {\small $-g\,c$} (N.south);
\draw[->,thick,blue!60!black] (-1.7,0) node[left,black] {\small $b$} -- (N);
\end{tikzpicture}
\caption{عصبون \nt{SERO} بتغذية راجعة عبر ناقله نفسه: يذهب التركيز $c$ إلى كل المتلقين ويثبّط العصبون نفسه. في النموذج هذا منظِّم مستوى؛ أما أن للحلقة هذا الدور في الدماغ ففرضية.}
\label{fig:serofeedback}
\end{figure}

لنحسب ما تفعله هذه الحلقة. لتكن خاصية النقل خطية، $f(u)=au$ عند $u>0$. في التوازن $c=\kappa r$ و$r=a(b-gc)$. وبتعويض إحداهما في الأخرى نحصل على
\begin{empheq}[box=\keybox]{equation}
r \;=\; \frac{a\,b}{1+G}, \qquad G \;=\; a\,g\,\kappa .
\label{eq:feedback}
\end{empheq}
المقدار $G$ هو كسب الحلقة. إنها الصيغة نفسها لأي مضخّم بتغذية راجعة سالبة، وتعطي المكسبين نفسيهما.

\emph{المناعة من التشتت.} لنفترض أن استثارية العصبون $a$ تغيّرت بنسبة $\varepsilon$. من دون تغذية راجعة كان التردد سيتغير بالنسبة نفسها. أما مع التغذية الراجعة، عند $G\gg1$، فيكاد التردد $r\approx b/(g\kappa)$ لا يعتمد على $a$: تغيّره أصغر بـ$1+G$ مرة. عند $G=9$ يُكبَت التشتت عشر مرات.

\emph{السرعة.} نعوّض $r=a(b-gc)$ في المعادلة الأولى من~\eqref{eq:seroneuron} فنحصل على $\tau_c\,\dd c/\dd t=-(1+G)\,c+\kappa ab$. يقترب التركيز من مستواه بثابت زمني $\tau_c/(1+G)$، أي أسرع بـ$1+G$ مرة مما هو من دون تغذية راجعة.

هذا أول حساب يمكن أن يؤديه نظام \nt{SERO}: أن يُبقي المستوى العام للسيروتونين في الدماغ عند قيمة مضبوطة، كما يُبقي الترموستات درجة الحرارة. هذه فرضية: في التجارب لا يظهر التثبيط الذاتي حتى الآن إلا توقفاتٍ قصيرة لا تغيّر التردد المتوسط.

\section{من النبضات إلى المستوى}

الحساب الثاني مخبوء في المعادلة الأولى من~\eqref{eq:seroneuron}. تطلق العصبونات نبضات منفردة، أما المتلقي فيحسّ بالتركيز، الذي ينعّمها على الزمن $\tau_c$. وإذا تغيّر تردد المصدر تبعه التركيز متأخرًا:
\begin{equation}
c(t) \;=\; \frac{\kappa}{\tau_c}\int_{-\infty}^{t} e^{-(t-t')/\tau_c}\,\sum_j r_j(t')\,\dd t' .
\label{eq:lowpass}
\end{equation}
هذا متوسط متحرك لنشاط المصادر على آخر بضعة $\tau_c$، أي مرشّح تمرير منخفض (الشكل~\ref{fig:lowpass}). وبهذه الطريقة يمرّر أبسط محوّل رقمي-تماثلي النبضاتِ عبر دارة \foreignlanguage{english}{RC} ويحوّل ترددها إلى مستوى. ونظام \nt{SERO} يفعل الشيء نفسه بمئات الآلاف من العصبونات دفعة واحدة.

وهذا المقدار ذاكرة أيضًا. بعد أن تصمت المصادر يبقى التركيز قائمًا زمنًا من رتبة $\tau_c$. إنه ليس بتًا، بل أثر يذوب تدريجيًا.

\begin{figure}[t]
\centering
\begin{tikzpicture}[>=Stealth,x=1.45cm,y=1cm]
\draw[gray!70!black] (0,3.2) -- (8.1,3.2);
\node[left,font=\small] at (-0.05,3.4) {النبضات};
\node[above,font=\scriptsize,gray!40!black] at (1,3.65) {$2$ في الثانية};
\node[above,font=\scriptsize,gray!40!black] at (3.5,3.65) {$8$ في الثانية};
\node[above,font=\scriptsize,gray!40!black] at (6.5,3.65) {$2$ في الثانية};
\draw[->,gray!70!black] (0,0) -- (8.3,0) node[right,black] {\small $t$};
\draw[->,gray!70!black] (0,0) -- (0,2.75) node[left,black] {\small $c$};
\foreach \s in {0.136,0.529,0.760,1.223,1.714,1.748,1.755,2.663,2.701,2.734,3.414,3.493,3.719,3.800,3.928,3.948,4.074,4.327,4.420,4.589,4.728,4.736,4.914,5.025,5.205,5.220,6.224,6.544,7.178}{\draw[thick,blue!60!black] (\s,3.2) -- (\s,3.6);}
\foreach \a/\b/\v in {0.000/0.136/0.000,0.136/0.529/1.000,0.529/0.760/1.675,0.760/1.223/2.330,1.223/1.714/2.466,1.714/1.748/2.509,1.748/1.755/3.425,1.755/2.663/4.402,2.663/2.701/2.775,2.701/2.734/3.672,2.734/3.414/4.553,3.414/3.493/3.306,3.493/3.719/4.055,3.719/3.800/4.235,3.800/3.928/4.905,3.928/3.948/5.316,3.948/4.074/6.211,4.074/4.327/6.476,4.327/4.420/6.028,4.420/4.589/6.493,4.589/4.728/6.483,4.728/4.736/6.642,4.736/4.914/7.589,4.914/5.025/7.351,5.025/5.205/7.579,5.205/5.220/7.331,5.220/6.224/8.221,6.224/6.544/4.012,6.544/7.178/3.914,7.178/8.000/3.076}{\draw[very thick,orange!85!black] plot[domain=\a:\b,samples=10] (\x,{0.25*\v*exp(-(\x-\a))});}
\foreach \s/\p/\q in {0.136/0.000/1.000,0.529/0.675/1.675,0.760/1.330/2.330,1.223/1.466/2.466,1.714/1.509/2.509,1.748/2.425/3.425,1.755/3.402/4.402,2.663/1.775/2.775,2.701/2.672/3.672,2.734/3.553/4.553,3.414/2.306/3.306,3.493/3.055/4.055,3.719/3.235/4.235,3.800/3.905/4.905,3.928/4.316/5.316,3.948/5.211/6.211,4.074/5.476/6.476,4.327/5.028/6.028,4.420/5.493/6.493,4.589/5.483/6.483,4.728/5.642/6.642,4.736/6.589/7.589,4.914/6.351/7.351,5.025/6.579/7.579,5.205/6.331/7.331,5.220/7.221/8.221,6.224/3.012/4.012,6.544/2.914/3.914,7.178/2.076/3.076}{\draw[very thick,orange!85!black] (\s,0.25*\p) -- (\s,0.25*\q);}
\draw[thick,densely dashed,black] plot[domain=0:2,samples=30] (\x,{0.25*2*(1-exp(-\x))})
  plot[domain=2:5,samples=40] (\x,{0.25*(8-6.271*exp(-(\x-2)))})
  plot[domain=5:8,samples=40] (\x,{0.25*(2+5.688*exp(-(\x-5)))});
\foreach \x in {0,2,5,8}{\draw[gray!60!black] (\x,-0.05) -- (\x,0.05) node[below=3pt,font=\scriptsize] {$\x\,\text{s}$};}
\end{tikzpicture}
\caption{من النبضات إلى المستوى. في الأعلى نبضات عصبونات \nt{SERO}: متفرقة ثانيتين، وكثيفة ثلاث ثوانٍ، ثم متفرقة من جديد. في الأسفل تركيز السيروتونين~\eqref{eq:lowpass} عند $\tau_c=1\,\text{s}$. الخط المتقطع هو التركيز نفسه لو جاءت النبضات بانتظام تام.}
\label{fig:lowpass}
\end{figure}

\section{السلك هو الحجم}
\label{sec:volume}

لنعد إلى شرط القضية~\ref{prop:serocomplete}. السيروتونين الذي يحرّره مرسِلون مختلفون متجاورون يتجمع في تركيز واحد.

\emph{السلك هنا ليس ليفًا، بل حجم من النسيج.} لا تبقى إشارتان متمايزتين إلا إذا حُرّرتا في منطقتين تفصل بينهما مسافة أكبر مما يتسع الوقت لانتشار الناقل فيه. خلال الزمن $\tau$ تبتعد جزيئة معامل انتشارها $D$ مسافة
\begin{empheq}[box=\keybox]{equation}
\ell \;\sim\; \sqrt{D\,\tau}\, .
\label{eq:diffusionlength}
\end{empheq}
يُبطئ تعرّجُ الشقوق بين الخلايا انتشارَ الجزيئة الصغيرة نحو $2.6$ مرة~\cite{Sykova2008}، أما معامل انتشار السيروتونين نفسه في النسيج فلم يُقَس؛ ومن حجم الجزيئة نقدّره ببضع وحدات من $10^{-6}\,\text{cm}^2/\text{s}$. إذا عاشت الإشارة $\tau\sim1\,\text{s}$ (وهذا تقدير أيضًا)، فإن $\ell\sim\sqrt{3\cdot10^{-6}}\,\text{cm}\approx20\um$. العنوان في شبكة كهذه هو \emph{المكان}: لا يعرف المتلقي مصدر الإشارة إلا من الموضع الذي يوجد هو فيه.

عصبونات \nt{SERO} الحقيقية مبنية بحيث لا يكاد يبقى شيء من هذه العناوين المنفصلة. خرج كل منها موزع على مساحات هائلة من الدماغ: فروع الخلية الواحدة تبلغ مناطق كثيرة متباعدة~\cite{GagnonParent2014}. يُقدَّر عدد مواقع التحرير في الخلية بنحو $10^5$ (الجدول~\ref{tab:types})، ولم يعدّها أحد مباشرة. ”أسلاك“ عصبونات \nt{SERO} المختلفة تتداخل وتختلط. ومع ذلك لا تندمج كليًا في سلك واحد: تنقسم عصبونات \nt{SERO} إلى بضعة أنظمة فرعية ترسل خرجها إلى مناطق مختلفة وتستجيب للحدث نفسه استجابات مختلفة~\cite{Ren2018}. لكن هذه الأنظمة الفرعية آحاد، لا مئات آلاف.

\begin{keyblock}
\begin{proposition}[عدد الأسلاك المستقلة]
\label{prop:volumewires}
في شبكة ذات نقل حجمي، لا يحدّ عددَ الإشارات المستقلة عددُ العصبونات، بل عددُ المناطق غير المتداخلة ذات الحجم $\ell\sim\sqrt{D\tau}$ التي تحرّر فيها الناقل. وإذا كان خرج كل عصبون موزعًا على حجم كبير، فلا تنقل الشبكة كلها إلا بضعة متغيرات مشتركة.
\end{proposition}
\end{keyblock}

لذلك لا يوجد في الدماغ ما تُبنى منه الدارات المنطقية للقضية~\ref{prop:serocomplete}. أما المنظِّم~\eqref{eq:feedback} والمرشّح~\eqref{eq:lowpass} فلا يحتاجان إلى أسلاك منفصلة: يكفيهما مقدار مشترك واحد. المرشّح ينتج مباشرة من التحرير في الحجم، وهو مقيس في مناطق الهدف~\cite{Bunin1998}؛ أما منظِّم المستوى ففرضية.

\section{أفق التخطيط}
\label{sec:horizon}

فيمَ قد ينفع مقدار مشترك بطيء واحد؟ أنسب ما يخطر هنا معامِل يجب أن يكون واحدًا للشبكة كلها، وألا يتغير أسرع من ثوانٍ أو دقائق. في الفصل~\ref{sec:neurontypes} ظهر معامل كهذا: المعامل $\gamma$ في صيغة الخطأ~\eqref{eq:rpe}، الذي يقول بكم تقلّ قيمة المكافأة المقبلة عن المكافأة الفورية. في الزمن المتصل يحل محله \emph{الأفق} $\tau_\gamma$: المكافأة $R$ التي يجب انتظارها زمنًا $D$ تساوي الآن $R\,e^{-D/\tau_\gamma}$.

الأفق يقرر هل يستحق الأمر الانتظار. لنفترض أنه يمكن أخذ مكافأة صغيرة $r_0$ فورًا أو انتظار مكافأة كبيرة $R$. يكون الانتظار مجديًا ما دام
\begin{empheq}[box=\keybox]{equation}
D \;<\; \tau_\gamma\,\ln\frac{R}{r_0} .
\label{eq:patience}
\end{empheq}
عند $\tau_\gamma=2\,\text{s}$ ومكافأة مؤجلة أكبر بثلاث مرات يستحق الأمر الانتظار حتى $2.2\,\text{s}$؛ وإذا كان الأفق أطول بمرتين، فحتى $4.4\,\text{s}$. الصبر يتناسب مع الأفق.

تستند الصيغة~\eqref{eq:patience} إلى خصم أسّي. أما الخصم المقيس على تأخيرات من رتبة الثواني فأقرب إلى المنحنى الزائدي $R/(1+D/\tau_\gamma)$: هكذا تنخفض عند القردة قيمةُ المكافأة المؤجلة في الاختيار، واستجابةُ عصبونات \nt{DOPA} للإشارة التي تُنبئ بها~\cite{KobayashiSchultz2008}. ومع المنحنى الزائدي يحل محل الشرط~\eqref{eq:patience} الشرطُ $D<\tau_\gamma\,(R/r_0-1)$: لم يعد الاعتماد على نسبة المكافأتين لوغاريتميًا بل خطيًا، ومع الأرقام نفسها تنزاح حدود الانتظار إلى الضعف تقريبًا. ويُستخرج المنحنى الزائدي من الأسّية نفسها إذا كان التأخير عشوائيًا (التمرين~\ref{ex:05:7})، ويبقى الصبر متناسبًا مع مقياس الزمن.

وفق فرضيةٍ ما، \nt{SERO} هو الذي يضبط الأفق~\cite{Doya2002}. ولديه كل ما يلزم لهذا الدور: مستوى مشترك واحد للشبكة كلها، يتغير خلال ثوانٍ. وثمة تجارب مباشرة أيضًا: إذا شُغّلت عصبونات السيروتونين اصطناعيًا، انتظر الحيوان المكافأة المؤجلة مدة أطول قبل أن يكفّ عن الانتظار~\cite{Miyazaki2014}، وواصل محاولة الحصول على المكافأة مدة أطول في المكان الذي صارت فيه نادرة~\cite{Lottem2018}. أما كيف يتحول تركيز السيروتونين إلى $\tau_\gamma$ داخل الشبكة فغير معروف. في الفصل التالي سيدخل الأفق في الخطأ الذي يبثّه \nt{DOPA}.

\section{الخلاصة}

\begin{table}[t]
\centering
\small
\begin{tabular}{>{\raggedright}p{3.6cm}>{\raggedright}p{4.3cm}>{\raggedright\arraybackslash}p{4.3cm}}
\toprule
& \nt{GLUT} & \nt{SERO} \\
\midrule
زمن الاستجابة & ميلي ثوانٍ & من أجزاء الثانية إلى ثانية \\
إشارة الأثر & $+$ فقط & $+$ و$-$، يختارها المتلقي \\
من دون دخل & يصمت & يعمل وحده مع دفق خلفي ثابت \\
العنونة & من نقطة إلى نقطة & حجم، والعنوان هو المكان \\
\foreignlanguage{english}{NOT} & عبر مستشعرات ”الإطفاء“ فقط & عصبون واحد \\
الذاكرة & نشاط ذاتي الاستمرار، يخبو بالإجهاد & تركيز، يذوب خلال $\tau_c$ \\
الحسابات الرئيسية & التزامن، والتأخيرات، والمنطق، والتسلسلات & التنعيم؛ وضبط المستوى والأفق $\tau_\gamma$ (فرضيتان) \\
\bottomrule
\end{tabular}
\caption{ما تحسبه شبكات من عصبونات \nt{GLUT} وحدها ومن عصبونات \nt{SERO} وحدها.}
\label{tab:glutvssero}
\end{table}

بوصفه عنصرًا منطقيًا (الجدول~\ref{tab:glutvssero}) عصبون \nt{SERO} أغنى من \nt{GLUT}، لكنه بوصفه نظام اتصال بثٌّ عام: بطيء ويكاد يخلو من العناوين. لذلك لا يحاول أن يكون معالجًا، بل ينعّم ويبثّ بضعة مقادير مشتركة تحدثنا عنها في الفصل~\ref{sec:neurontypes}، ومنها، وفق فرضيةٍ، أفق التخطيط. وسنلتقي ناظمة الإيقاع والحجم والتركيز البطيء ثلاث مرات أخرى: عند \nt{DOPA} و\nt{NORA} و\nt{ACET}.

\section{تمارين}

\begin{exercise}[\lvlA]
\label{ex:05:1}
\emph{السلك هو الحجم.} خذ للسيروتونين $D=3\cdot10^{-6}\,\text{cm}^2/\text{s}$ (القسم~\ref{sec:volume}). جد من~\eqref{eq:diffusionlength} قيمة $\ell$ لإشارة تعيش $1\,\text{s}$، والزمن الذي ينتشر فيه الناقل مسافة $5\,\text{cm}$. فبمَ إذن يتحقق البثّ في \nt{SERO}: بالانتشار أم بتفرّع سلك له $\sim10^5$ موقع تحرير؟
\end{exercise}

\begin{exercise}[\lvlA]
\label{ex:05:2}
\emph{منظِّم المستوى.} بيّن أنه في~\eqref{eq:seroneuron} مع $f(u)=au$ تساوي الحساسية النسبية $S=(a/r)\,\partial r/\partial a$ القيمةَ $1/(1+G)$ تمامًا، لا عند $G\gg1$ فقط. عند $G=9$ ازدادت الاستثارية $a$ بنسبة $20\,\%$. بكم في المئة يتغير $r$ من دون تقريب خطي؟
\end{exercise}

\begin{exercise}[\lvlB]
\label{ex:05:3}
\emph{مستقبل بطيء في الحلقة.} يستجيب مستقبل \nt{SERO} بتأخير: $r(t)=a\bigl(b-gc(t-T)\bigr)$، والحجم وفق~\eqref{eq:seroneuron}. بيّن أنه عند $G>1$ لا تكون الحلقة مستقرة إلا إذا كان $T<T^*=\tau_c\arccos(-1/G)/\sqrt{G^2-1}$. جد $T^*$ عند $\tau_c=1\,\text{s}$ لـ$G=2$ و$G=9$. ما كبت التشتت المتاح مع تأخير من مئات الميلي ثانية؟
\end{exercise}

\begin{exercise}[\lvlB]
\label{ex:05:4}
\emph{ضجيج الطلقات في المستوى.} كل نبضة تضيف إلى $c$ وفق~\eqref{eq:lowpass} دالة أسّية بـ$\tau_c=1\,\text{s}$. جد معامل الاختلاف لـ$c$ إذا كانت كل عصبونات \nt{SERO}، وعددها $3\cdot10^5$، تحرّر في حجم واحد نبضات بواسونية بمعدل $2$ في الثانية. أي $\tau_c$ يعطي $\mathrm{CV}=10\,\%$ لعصبون واحد تردده $4$ نبضات في الثانية؟
\end{exercise}

\begin{exercise}[\lvlB]
\label{ex:05:5}
\emph{نمذجة: التغذية الراجعة ضد الضجيج.} حاكِ $N=50$ ناظمة إيقاع بواسونية بترددات $r_i=a_i[b-gc]_+$، حيث $a_i$ موزعة بانتظام في $[2.8,\,5.2]$، والحجمَ~\eqref{eq:seroneuron} مع $\kappa=1/N$ و$\tau_c=1\,\text{s}$. كم مرة تخفّض التغذية الراجعة ($b=1$، $g=2.25$) معامل الاختلاف \foreignlanguage{english}{CV} للمستوى $c$ مقارنةً بـ$b=0.1$، $g=0$؟ وهل تسوّي ترددات العصبونات؟
\end{exercise}

\begin{exercise}[\lvlB]
\label{ex:05:6}
\emph{تصميم: \foreignlanguage{english}{XOR} بمنطق \nt{SERO}.} البوابة ناظمة إيقاع تعطي $r_0$ إذا كان $b-g\sum_kc_k>0$، و$0$ خلاف ذلك، في حجمها الخاص $\tau_c\,\dd c/\dd t=-c+\kappa r$؛ وهي تحسب \foreignlanguage{english}{NOR}. ابنِ \foreignlanguage{english}{XOR} من خمس بوابات كهذه، وجد زمن استقراره عند $\tau_c=1\,\text{s}$ و$G'=g\kappa r_0/b=2$.
\end{exercise}

\begin{exercise}[\lvlB]
\label{ex:05:7}
\emph{الصبر مع تأخير مجهول.} (أ)~يقبل الحيوان انتظار مكافأة أكبر بثلاث مرات إذا لم يتجاوز الانتظار $6\,\text{s}$. أي أفق $\tau_\gamma$ يعني هذا؟ (ب)~ليكن التأخير $D$ موزعًا أسّيًا بمتوسط $\bar D$. بيّن أن الانتظار مجدٍ عند $\bar D<\tau_\gamma(R/r_0-1)$، وقارن بالتأخير الثابت عند $\tau_\gamma=2\,\text{s}$ و$R=3r_0$.
\end{exercise}

\begin{exercise}[\lvlC]
\label{ex:05:8}
\emph{كيف يضبط التركيز الأفق.} يتلقى عصبون \nt{DOPA} المقدار $V$ مباشرة بوزن $k_1$، وعبر وسيط \nt{GABA} ينعّمه بـ$\tau_d=0.1\,\text{s}$ بوزن $-k_2$؛ وللمقادير $V$ البطيئة يساوي المجموع $(k_1-k_2)V+k_2\tau_d\dot V$. اختر $k_1$ و$k_2$ لتطابقا~\eqref{eq:ctd} مع $\tau_\gamma=2\,\text{s}$. يضرب \nt{SERO} المعامل $k_1$ في $m$: كم يجب أن يتغير $m$ لمضاعفة الأفق؟
\end{exercise}
